% -------------------------------------------------
% VERSION TRACKER — No modificar este bloque
% Versión:    Tesis Humanizada
% Archivo:    Capitulo1Introduccion/Formulacion.tex
% Última mod:  2026-05-08T08:16:19
% Git commit:  cc9c766f @ main
% Audiencia:   general
% Verificar antes de editar: bash check_tesis_version.sh overleaf_tesis_humanizado/Capitulo1Introduccion/Formulacion.tex
% -------------------------------------------------


\section{Formulación del problema}

¿Cómo desarrollar un sistema web que, además de digitalizar el control de inventarios, incorpore un asistente inteligente capaz de responder preguntas en lenguaje natural, para ayudar a las PYMEs colombianas a reducir pérdidas por vencimiento, ganar trazabilidad en sus movimientos de bodega y tomar decisiones basadas en datos reales —dejando atrás la dependencia del cuaderno y la memoria?

\section{Descripción del problema}

Las pequeñas y medianas empresas (PYMEs) son el motor de la economía colombiana: representan el 99,5\% de las empresas formales del país y generan cerca del 40\% del Producto Interno Bruto (PIB).\cite{bbva2024mipymes} Sin embargo, existe una brecha importante en su nivel de digitalización. Mientras que el 63\% de las empresas medianas ya ha avanzado en su transformación digital, solo el 42\% de las microempresas ha dado ese paso.\cite{innpulsa_transformacion}

% El corazón del problema: la gestión manual del inventario
Una de las áreas más rezagadas es el control de la bodega. Muchas PYMEs todavía registran sus entradas y salidas de materia prima en cuadernos, libros de contabilidad o, en el mejor de los casos, hojas de cálculo que cada quien maneja a su manera.\cite{bbva2024mipymes} Estos métodos manuales tienen un costo real: errores de conteo, registros duplicados o perdidos, y una incapacidad casi total de saber qué pasó con un insumo, cuándo y por qué. La situación es más grave si se considera que apenas el 40\% de las microempresas colombianas usa internet en sus operaciones diarias.\cite{bbva2024mipymes}

% Las consecuencias, en términos concretos
Esta forma de trabajar produce tres problemas que golpean directamente el bolsillo de la empresa \cite{ballou2004logistica}:

\begin{itemize}
    \item \textbf{Pérdidas por ineficiencia}: Productos que se vencen o se dañan antes de usarse, falta de insumos justo cuando más se necesitan y compras de emergencia a precios elevados porque nadie anticipó la escasez.
    \item \textbf{Falta de visibilidad}: El dueño o administrador no sabe en tiempo real qué tiene en su bodega. Sin esa información, planificar compras se convierte en un juego de adivinanza.
    \item \textbf{Decisiones a ciegas}: Sin datos históricos ni reportes confiables, las decisiones estratégicas (¿cuánto comprar?, ¿cada cuánto?, ¿a qué proveedor?) se toman por intuición y experiencia, no por evidencia.
\end{itemize}

Aunque existen sistemas de gestión de inventarios en el mercado —como Odoo, Zoho Inventory o Stockpile— la mayoría resultan inalcanzables para una PYME colombiana. Sus precios de suscripción, la complejidad de sus interfaces y la necesidad de personal técnico para operarlos los convierten en herramientas diseñadas para grandes empresas, no para panaderías, ferreterías o talleres.\cite{bbva2024mipymes}

\subsection{Caso de estudio: una panificadora del Valle del Cauca}

Para entender el problema desde la operación real y no solo desde la literatura, se realizó un levantamiento de requisitos con una panificadora de tamaño pequeño a mediano del municipio de La Victoria (Valle del Cauca), que elabora productos de panadería a escala industrial y compra de forma continua harinas, mejoradores, levaduras, grasas, semillas, sal, azúcar y material de empaque. Por respeto a la empresa y a las personas involucradas, en este documento no se registran sus nombres.

El levantamiento consistió en una entrevista semiestructurada de aproximadamente 41 minutos con la persona a cargo de la administración de la empresa, organizada en diez preguntas en lenguaje cotidiano y un grupo de preguntas técnicas depuradas. El detalle completo de la entrevista y de los requerimientos que se derivaron de ella se presenta en el Capítulo \ref{chap:requerimientos}.

La entrevista mostró que el problema de esta empresa no es la pérdida de materia prima por vencimiento ---la rotación es rápida y, según la administradora, no se presentan mermas por caducidad--- sino la calidad de la información con la que se gestiona la bodega:

\begin{enumerate}
    \item \textbf{El inventario digital no coincide con la bodega}: El sistema contable se alimenta a partir de las facturas y del formato diario de producción, sin verificación física. La confusión entre referencias casi idénticas al digitar las facturas y los consumos que no se desglosan generan diferencias que solo aparecen en el cierre de mes.
    \item \textbf{La conciliación es manual}: Una vez al mes, la administradora compara en una hoja de cálculo el conteo físico con el saldo del sistema contable y corrige las diferencias con movimientos manuales que no dejan constancia de quién los hizo ni por qué.
    \item \textbf{El reabastecimiento depende de la experiencia de una persona}: No hay stock mínimo ni punto de reorden; los pedidos se deciden recorriendo la bodega. La empresa ya ha tenido faltantes de insumos que obligaron a comprarlos en el mercado local a un precio mayor.
    \item \textbf{La trazabilidad por lote es una exigencia sanitaria}: Aunque no haya mermas, ante una inspección del INVIMA la empresa debe demostrar qué lotes de materia prima usó y cuándo vencían.
\end{enumerate}

Estos problemas no son exclusivos de esta empresa: coinciden con las dificultades de digitalización que las fuentes citadas describen para las PYMEs colombianas, y confirman la pertinencia de una herramienta de gestión de inventario orientada a este segmento.

\subsection{La brecha entre las PYMEs y el software disponible}

La Tabla \ref{tab:barreras_pymes} resume las barreras que separan a las PYMEs colombianas de las soluciones de inventario que existen hoy en el mercado.

\begin{table}[H]
\centering
\caption{Barreras para la adopción de software de inventario en PYMEs colombianas.}
\label{tab:barreras_pymes}
\begin{tabular}{|l|p{9cm}|}
\hline
\textbf{Barrera} & \textbf{Descripción} \\
\hline
Costo de licenciamiento & Suscripciones mensuales desde \$200.000 COP hasta \$2.000.000 COP por usuario \\
\hline
Complejidad técnica & Interfaces diseñadas para grandes empresas con procesos estandarizados \\
\hline
Falta de contexto local & Software diseñado en otros países sin adaptación al español colombiano ni a las regulaciones locales \\
\hline
Curva de aprendizaje & Tiempo y costo de capacitar al personal en el uso de la herramienta \\
\hline
Dependencia de infraestructura & Requisitos de hardware, servidores o conectividad permanente \\
\hline
\end{tabular}
\end{table}

A estas barreras se suma una ausencia llamativa: ninguna de las soluciones disponibles ofrece un módulo que permita preguntar en lenguaje natural cosas como ``¿qué debo despachar primero?'' o ``¿cuánta harina me queda?'' sin tener que navegar por menús complejos o aprender filtros avanzados. El operario de bodega necesita respuestas, no capacitación en software.

\subsection{Qué aporta Pymetory frente a este problema}

Frente al panorama descrito, Pymetory se diseñó para cerrar cinco brechas concretas:

\begin{itemize}
    \item \textbf{Kardex digital automático}: Reemplaza los registros manuales por un sistema que documenta cada movimiento automáticamente, guardando quién lo hizo, cuándo, de qué tipo fue, qué cantidad movió y por qué se realizó.
    \item \textbf{Despacho inteligente (FEFO)}: El sistema toma siempre primero el lote con la fecha de vencimiento más próxima, formalizando la práctica de ``gastar primero lo que ya estaba en bodega'' que hoy depende de la memoria del personal.
    \item \textbf{Consultas en español colombiano}: El usuario puede preguntar ``¿cuánta harina me queda?'', ``¿qué está por vencer?'' o ``¿quién hizo el último despacho?'' y recibir una respuesta basada en los datos reales de \textit{su} inventario.
    \item \textbf{Alertas automáticas}: El sistema avisa cuando un lote está próximo a vencer o cuando el stock de un material cae por debajo del nivel mínimo, sin depender de que alguien se acuerde de revisar.
    \item \textbf{Reportes en un clic}: Informes de inventario, movimientos, FEFO, consumo y valorización en formatos PDF, CSV y Excel, que apoyan el cierre y la conciliación del inventario.
\end{itemize}

\newpage
